\section{Leis estatísticas e Sistemas complexos}


Altmann  nos oferece uma abordagem unificadora sobre o uso das leis estatísticas na investigação dos sistemas complexos, isto é, existe uma abordagem única que é aplicada na investigação de diferentes leis estatísticas\cite{Altmann_2024}. Esse padrão na forma como a investigação é conduzida para diferentes casos particulares é um dos primeiros sinais  de que esse conjunto de casos particulares constitui uma área de investigação dotada de uma metodologia científica particular. Veremos mais adiante como também utilizamos esta abordagem na investigação da econofísica.

Começaremos,  então, reconhecendo que a identificação de padrões é uma etapa de grande importância no processo de investigação científica de forma abrangente. Esta é uma etapa fundamental da investigação científica na própria física, como é tradicionalmente concebida. Dessa forma, a primeira conexão que podemos estabelecer entre a econofísica e a física tradicional consiste exatamente na centralidade do processo de identificação de padrões na prática  científica.  Uma das primeiras exigências da econofísica consiste, então, na aplicação desta ferramenta científica mais geral, evidentemente agora aplicada dentro do contexto econômico; ou seja, isto nada mais é que uma ``extensão'' de uma prática científica já consolidada.  

Porém, esta afirmação deve ser realizada com certo cuidado, uma vez que o estudo das leis estatísticas remonta ao nascimento de muitas disciplinas científicas no século XVII. A origem da ideia de que diferentes conjuntos de dados e fenômenos podem ser descritos por uma mesma função ou distribuição universal está intimamente relacionada com à tentativa de alargar o sucesso dos métodos quantitativos nas ciências naturais (especialmente a física clássica) para as ciências biológicas e sociais por meio da da estatística. Esta ideia desempenha um papel central nos trabalhos do cientista francês Pierre-Simon Laplace (1749-1827) e do belga Adolphe Quételet (1796-1874) na primeira metade do século XIX. Quételet, por exemplo, buscou identificar padrões em dados relativos ao nascimento, à idade de casamento, a atividades criminosas e às taxas de mortalidade.

Karl Marx (1818-1883), no século XIX, também propunha uma aproximação entre o que entendemos hoje como as áreas gerais das ciências da natureza e das ciências humanas. \footnote{``Conhecemos uma única ciência, a ciência da história. A história pode ser examinada de dois lados, dividida em história da natureza e história dos homens. Os dois lados não podem, no entanto, ser separados;'' \cite{marx2021ideologia}}. Companheiro intelectual de Marx,  Friedrich Engels (1820-1895) segue pelo mesmo caminho e, de forma mais explícita, sugere a existência de leis internas universais no desenvolvimento da história humana, análogas às leis naturais, de certa forma, adiantando que a multiplicidade de interesses e objetivos de cada indivíduo que compõe um sistema exige uma modelagem que seja capaz de lidar com o `acaso'\footnote{``... o curso da história é regido por leis internas universais. Pois, também aqui, apesar dos objectivos conscientemente queridos de todos os indivíduos, domina aparentemente à superfície, grosso modo, o acaso. Só raramente acontece o querido; na maioria dos casos, os múltiplos objectivos queridos entrecruzam-se e contradizem-se, ou esses mesmos objectivos são de antemão irrealizáveis, ou os meios são insuficientes. Assim, os choques das inúmeras vontades individuais e acções individuais conduzem a um estado que é totalmente análogo ao que domina na Natureza desprovida de consciência.''\cite{engels2024ludwig}.}.

O termo ``física social'' ficou associado, então, a este campo emergente de estudos sociais quantitativos. Este termo também foi adotado pelo positivista francês Auguste Comte (1798-1857). Embora Comte tenha mais tarde transitado para os termos ``Sociologia'' e ``Ciências Sociais'', que se tornaram mais prevalentes, a ``sociofísica'' ou ``física social'' persistiu no século XX e continua a ser utilizada, frequentemente associada a modelos inspirados na física (da matéria condensada). 

Antes de prosseguirmos, faz-se necessária uma definição de sistemas complexos: sistemas complexos são sistemas compostos por vários componentes microscópicos que interagem entre si, dando origem a fenômenos não triviais em escalas macroscópicas. Esses fenômenos, que não podem ser encontrados nos componentes individuais do sistema, são chamados de fenômenos emergentes, isto é, propriedades que emergem a partir da interação entre os componentes do sistema.

Neste contexto de sistemas complexos, esses padrões costumam ser descritos não por meio de leis determinísticas, como podemos encontrar, por exemplo, na mecânica clássica, mas sim por meio de leis estatísticas.  Precisamos, então, definir o que queremos dizer com lei estatística: no contexto dos sistemas complexos, as leis estatísticas são propriedades emergentes com características estatísticas inerentes. A explicação dessas leis deve ser obtida por meio da consideração de modelos microscópicos que conduzem à observação da lei estatística em escalas macroscópicas.

Podemos então presumir que a lei estatística é uma lei matemática que diferentes sistemas exibirão; isto é necessário para fazer jus ao nome de ``lei''. Ela também é necessariamente de natureza estatística; isto é, a lei se dará em termos de grandezas estatísticas, o que é necessário para fazer jus ao nome de ``estatística''. Ela se dará, por exemplo, em termos de distribuições, probabilidades, médias, etc; ou seja, valores que só podem ser obtidos através de uma grande quantidade de medidas. E também temos que, por definição, esta lei é necessariamente uma propriedade emergente de um sistema complexo.

A partir dessas definições, somos capazes tanto de começar a delimitar qual é o nosso objeto de estudo (sistemas complexos que apresentam uma lei estatística como propriedade emergente) quanto de caracterizar o método de investigação adotado (a descrição de padrões observados no sistema por meio de leis estatísticas em um primeiro momento e a utilização de modelos microscópicos capazes de gerar essa lei macroscópica em um segundo momento). Evidentemente, a validade de uma lei estatística dependerá de qual sistema estamos analisando e do seu contexto; só podemos esperar que uma lei estatística sobre a distribuição de riqueza em países capitalistas  seja válida dentro do contexto dos países capitalistas. 

Altmann descreve esta forma de fazer ciência como um novo paradigma na investigação da realidade social \cite{Altmann_2024}. Para entendermos melhor as implicações desta afirmação, podemos trabalhar com um exemplo.  Por volta do fim do século XX, o físico Felix Auerbach (1856-1933) notou uma regularidade na população das cidades da Alemanha. Quando fez um ranking das cidades, tendo como critério a sua população, sendo  $r=1$  a cidade mais populosa e $r=R$ a menos populosa, e a população da cidade $r$  sendo  dada por $N_{r}$,  então foi obtida a seguinte lei estatística:

$$N_{r}=\frac{A}{r}$$

Onde $A$ é uma constante aproximadamente igual a população da cidade mais populosa $N_1$. Auerbach percebeu que a observação valia também para outros países e, através de generalizações posteriores, tornou-se uma das leis estatísticas mais famosas, a lei do tamanho das cidades de Auerbach-Lotka-Zipf.

O potencial das leis estatísticas pode ser visto neste exemplo. Elas são funções com poucos parâmetros ($A$) e são propostas para descrever um grande conjunto de dados (cidades) em diferentes cenários (países). Isso permite não apenas um resumo dos dados, mas também análises analíticas, além de abrir a possibilidade para reflexões teóricas. Os estudos das leis estatísticas não devem ser encarados como meras curiosidades; a investigação dessas leis pretende revelar as propriedades mais fundamentais do sistema investigado.

Na tradição da física social, as leis estatísticas eram vistas como semelhantes às leis empíricas da física. Sendo assim, o termo física social estava associado à busca de comportamentos semelhantes a leis na sociedade. A ideia de que existiam leis relacionadas à sociedade, assim como a mecânica de Newton se relacionava com o movimento dos planetas, era compartilhada por muitos cientistas no final do século XVII. A expectativa da física social era a de que ela evoluiria como disciplina de forma semelhante ao desenvolvimento historicamente construído da mecânica clássica: 
\begin{center}
\textbf{Observações $\rightarrow$ Leis Empíricas $\rightarrow$  Leis Universais.}
\end{center}

Ou seja, primeiro a comunidade científica realiza um grande conjunto de observações, prossegue, então, para a extração de diversas leis empíricas pontuais dessas observações e, como último passo, a partir das leis empíricas, abstrai leis mais gerais que explicam tanto as leis empíricas para diferentes casos particulares quanto os dados observados como um todo. Como exemplo de leis empíricas, podemos citar Kepler e Galileu, e como lei universal, temos Newton. Podemos entender o escopo como a diferença fundamental, enquanto as leis empíricas têm uma validade mais limitada uma lei universal nos fornece uma descrição mais fundamental e universal da realidade, incluindo a própria explicação de diversas leis empíricas.

Nessa perspectiva, as leis estatísticas desempenham o papel de leis empíricas, constituindo a etapa intermediária crucial entre as observações empíricas e o desenvolvimento de teorias de validade geral. Isto é, algo temporário, até que sejamos capazes de desenvolver uma lei universal. Porém, esta analogia simplista ignora a natureza estatística crucial das leis estatísticas, que são fundamentalmente diferentes da lei de Kepler. 

Até o momento, as expectativas para o desenvolvimento da física social seguindo este caminho não foram confirmadas, nem há indícios de que surgirá uma teoria unificada semelhante à Mecânica Clássica. No entanto, um aspecto útil dessa visão ingênua, mantido nas aplicações contemporâneas das leis estatísticas, é a expectativa de que elas conectem observações a modelos teóricos, mesmo que isto não proceda de forma totalmente análoga ao desenvolvimento da mecânica newtoniana. 

As primeiras leis estatísticas foram fundamentais para o nascimento da estatística como disciplina, uma vez que transmitiram o potencial do pensamento estatístico e probabilístico a todas as ciências. O desenvolvimento dessas ideias descreveu com sucesso muitas das regularidades observadas como sendo consequência de flutuações aleatórias e também contribuiu para o desenvolvimento da própria Física através dos trabalhos de Maxwell (1831-1879) e Boltzmann (1844-1906) na Mecânica Estatística. Uma reviravolta irônica dos acontecimentos, uma vez que a física surge como um modelo padrão de ciência, mas acaba sendo ela mesma revolucionada com o advento da mecânica estatística. 

As leis estatísticas da física social diferem, então, das leis empíricas  da física tradicional no sentido de que não necessariamente possuem um caráter ``transitório'', existindo apenas enquanto aguardamos o desenvolvimento de uma teoria universal mais geral da qual podemos deduzir a lei estatística. Porém, isso não impede que a lei estatística conecte a observação a um modelo teórico mais geral. 

Por exemplo, ainda que obtenhamos uma formulação microscópica satisfatória das regras de interação entre os agentes que leva a uma distribuição de renda que obedece a lei estatística de Pareto, a partir deste conjunto de leis não podemos obter diretamente a distribuição. O que podemos fazer é realizar uma simulação e, com os dados gerados a partir dessa simulação, podemos, então, negar (ou não) que este conjunto de regras é capaz de gerar a lei estatística investigada. Mas, se quiséssemos realizar uma análise analítica sobre a distribuição, possivelmente retomaríamos a lei estatística de Pareto. Neste exemplo, a lei de Pareto não pode ser deduzida diretamente a partir da descrição microscópica do sistema  e conserva um papel a desempenhar na investigação deste sistema. 

Além disso, tanto as leis estatísticas quanto os modelos microscópicos podem se inserir em uma discussão teórica mais ampla sobre o fenômeno analisado. Segundo alguns teóricos, a própria econofísica  pode ser vista como um braço da economia. Um outro desenvolvimento possível das pesquisas nesta área consiste na proposta não de uma lei universal, mas sim de um modelo  microscópico mais geral capaz de reproduzir não apenas uma, mas múltiplas leis estatísticas.

Dessa forma, é o uso combinado das leis estatísticas e modelos microscópicos nos sistemas complexos que caracteriza esta abordagem. Há três mudanças que merecem ser destacadas na forma como se trabalhava na física social, causadas pela adesão a abordagem baseada em sistemas complexos:

\begin{enumerate} 
    \item Não ver as leis estatísticas como um simples efeito independente e agregado de influências aleatórias que se espera encontrar em partes individuais do sistema, mas sim como uma propriedade emergente de um sistema complexo. 
    \item O interesse nas leis vai além de aplicações práticas ou discussões filosóficas; são motivações e justificativas para propostas de modelos mecanísticos dos fundamentos microscópicos do sistema. 
    \item A explicação teórica das leis não segue o paradigma da mecânica clássica, mas o da mecânica estatística, isto é, baseada em modelos probabilísticos. 
\end{enumerate}

Esta abordagem começou a se desenvolver na metade do século XX, e entre os principais nomes, podemos citar Benoit Mandelbrot. Para que a sociofísica se desenvolvesse como existe hoje, foi necessário tanto o desenvolvimento teórico na área de sistemas complexos quanto o desenvolvimento material da sociedade que permitisse a existência de uma forma acessível de um poder computacional cada vez maior. Cientistas do passado, ainda que estivessem munidos das melhores intenções, encontraram-se necessariamente limitados pelo contexto em que estavam. 

No mesmo sentido, a disponibilidade atual de grandes conjuntos de dados, diferente de períodos anteriores, é o que permite que, atualmente, confrontemos as leis estatísticas propostas com dados reais. Assim, somos capazes de estender suas aplicações a novos casos, obter novas leis ou abandonar leis que foram formuladas erroneamente, se for o caso. 


\subsection{Formalização}

Altmann nos fornece uma definição mais robusta de lei estatística\cite{Altmann_2024}:

\textbf{Definição}: uma lei estatística (em sistemas complexos) é uma função que:

\begin{enumerate} 
    \item Foi proposta para descrever um grande número de observações em diferentes cenários (universalidade); 
    \item  É elementar ou uma composição de funções elementares com um pequeno número de parâmetros e dimensões (simplicidade); 
    \item Desempenha um papel importante em uma teoria ou modelo (conexão teórica). 
\end{enumerate}

Tipicamente, uma lei estatística se aplica a um conjunto de dados e descreve a frequência de observação de alguma coisa ou a relação entre duas propriedades do sistema observado. É esperado que a lei seja válida para todos os casos similares, não se restringindo a apenas um caso. A segunda condição pode ser expressa de forma matemática, afirmando que a dimensão da função deve ser muito menor que o número de observações. Existe uma ideia implícita de que se deve buscar a menor quantidade possível de parâmetros, visando a simplicidade. E, por fim, estas leis possuem, então um papel central de servir como uma conexão teórica entre os dados observados e modelos microscópicos que devem fornecer uma explicação mecanicista da lei e/ou serem usados em uma teoria mais geral sobre os fundamentos deste sistema.

É interessante apresentarmos um contra-exemplo. A simples distribuição gaussiana da altura dos humanos não é uma lei estatística nesta definição, pois falha na condição (3). Algo semelhante pode ser afirmado sobre modelos construídos para serem ajustados a um dado específico. Também não são leis estatísticas, pois não possuem validade universal (1) e nem provém uma ideia que favoreça a intuição (\textit{insight})  sobre uma teoria mais geral que explique os dados (3).

É natural a existência de um certo ceticismo sobre a possibilidade de generalizar algo como ``leis estatísticas'' e mesmo sobre a existência de regularidades universais na economia \cite{GALLEGATI20061}.  Rickles reconhece que, por essa dúvida estar baseada na ideia de livre arbítrio do homem,  a defesa da existência de leis a níveis coletivos não é uma questão fácil de responder. No entanto, ele considera necessário também reconhecer que não somos livres para agir de literalmente qualquer maneira. O comportamento das pessoas é limitado pela condição social em que se encontram, assim como por restrições biológicas, como a natureza física do cérebro e do corpo humano. Assim, em algum nível, é possível aceitar à ideia da existência de leis que restringem a distribuição empírica de possíveis comportamentos humanos, ainda que de forma limitada e aberta a debate\cite{RICKLES2007948}.

Da mesma forma que é possível considerar justo o ceticismo quanto à existência de leis no nível coletivo, também podemos ser céticos quanto a ausência dessas leis. Afinal, como Rickles coloca, esta também é uma suposição sobre como o mundo é e que é apresentada sem argumentos de apoio que sejam mais fortes. De toda forma,  a ciência deve estar pronta para ser reexaminada no que diz respeito à existência das leis sempre que necessário.

A defesa da  existência da universalidade de leis estatísticas decorre da observação de que estas leis aparentam serem motivadas, justificadas, usadas e estudadas de forma similar. De maneira geral, a abordagem de raciocinar com as leis estatísticas segue uma sequência de três passos conectados\cite{Altmann_2024}:

\textbf{1. Definição e validação}: Começamos com uma análise empírica. O passo inicial é, inevitavelmente, análogo a toda ciência. A análise de um grande conjunto de observações. É neste momento que vamos observar algum padrão nos dados e propor uma lei estatística. Também é neste momento que vamos testar a validade da lei proposta em outros conjuntos de dados. Neste passo, também podem ocorrer diversas reformulações, generalizações, modificações e toda a sorte de processos científicos tradicionais. 

Na sequência, os próximos dois passos acontecem de forma paralela; isto é, em posse da lei estatística proposta no passo anterior, temos dois possíveis caminhos a percorrer:

\textbf{2a. Explicação da lei}: a explicação da lei estatística através de um modelo generativo, isto é, um modelo microscópico capaz de gerar uma distribuição que obedeça à lei estatística. Após ser considerado que a lei estatística é empiricamente válida, uma primeira possível questão a ser investigada é sua origem. É com isso em mente que propomos modelos mecanísticos simples que gerem observações que satisfaçam às leis estatísticas. 

\textbf{2b. Consequências da lei}: uma segunda pergunta que surge após considerarmos que a lei estatística é válida é quais são as consequências dessa lei. Uma extrapolação de suas consequências para além do que observamos até então, isto é, predições baseadas na própria lei, na sua relação com outras leis, nas consequências do uso desta lei como restrição em outros modelos, etc… A capacidade preditiva de uma lei é algo tipicamente esperado em disciplinas científicas, especialmente nas ciências da natureza

De modo sucinto, a investigação de sistemas complexos por meio de suas leis estatísticas é possível com a existência de três elementos fundamentais: o sistema complexo a ser investigado, uma lei estatística (macroscópica) emergente deste sistema e um modelo generativo (microscópico) capaz de reproduzir a lei estatística. Este é o método padrão de investigação explorado por Altmann e adotado nesta pesquisa.

\subsection{Modelo generativo microscópico}

Não há, por definição, uma única abordagem possível para a construção do modelo microscópico. Vamos adotar a modelagem baseada em agentes (MBA) como o conceito usado para construir nossos modelos. Sobre isso, uma recomendação de leitura é o livro publicado pelos desenvolvedores do NetLogo\cite{10.2307/j.ctt17kk851}, um popular software voltado para a modelagem baseada em agentes.

De forma geral, um modelo é sempre uma abstração da realidade, ou seja, uma representação simplificada que preserva uma equivalência adequada para determinadas situações, de forma que pode ser vista como uma visão substitutiva da realidade para facilitar o raciocínio %\apud{phdthesis}{goldberg}.

A eficiência de um modelo depende da correta tradução da realidade para o modelo no ato de modelar, que, por sua vez, é exatamente o ato de simular a realidade através de modelos. É por isso que é importante pensar no modelo antes do próprio ato de modelagem começar, identificar as variáveis relevantes para o problema que será abordado e definir a linguagem que será utilizada, pois a elaboração do modelo depende desses pontos cruciais \cite{phdthesis}.  Como um modelo é uma simplificação da realidade, é possível que uma mesma realidade seja modelada de diferentes formas, já que nenhum modelo é a realidade em si. Temos como exemplo a forma mais tradicional de modelagem científica de processos dinâmicos: equações diferenciais (\textit{equation based modeling} - EBM). 

Outra forma de modelagem que tem se tornado cada vez mais popular é a modelagem baseada em agentes (\textit{agent based model} - ABM). Esta é uma forma de modelagem computacional na qual a dinâmica de um sistema complexo (social, biológico, etc) é representada em termos dos agentes e suas interações. Os agentes, por suas vez, são indivíduos computacionais autônomos que possuem propriedades e regras de comportamento que regem sua interação com outros agentes e com o ambiente.

Conforme apresentamos anteriormente, sistemas complexos são definidos como sistemas compostos por múltiplos elementos individuais que interagem entre si e apresentam um comportamento que não pode ser previsto a partir dos próprios elementos \cite{10.2307/j.ctt17kk851}. Como resultado das interações entre os elementos distribuídos, surge, então, um ``fenômeno emergente'', característico de sistemas complexos. Por sua vez, o fenômeno emergente possui duas características essenciais: A primeira é que as estruturas que emergem não podem ser deduzidas apenas das propriedades dos elementos, pois surgem da interação entre eles; ou seja, são propriedades do sistema e não dos elementos que o compõem. A segunda característica é a auto-organização, isto é,  o fenômeno surge espontaneamente sem a presença de um coordenador central. A emergência pode ser compreendida como o processo que mantém esta estrutura.

A ideia central da modelagem baseada em agentes é a proposta de que a maioria dos fenômenos do mundo real podem ser efetivamente modelados com agentes, um ambiente e uma descrição das interações agente-agente e agente-ambiente. O objetivo deste tipo de modelagem é criar agentes e regras que gerem o comportamento desejado. Dentro da abordagem que estamos discutindo, o comportamento desejado é apresentado na forma de leis estatísticas. Uma vez criados, esses sistemas podem ser usados para obter um melhor entendimento do fenômeno modelado por meio da experimentação de diferentes regras e propriedades, visando testar diferentes hipóteses.

\cite{10.2307/j.ctt17kk851} cita 8 principais usos para um modelo baseado em agentes:
\begin{itemize}
    \item Descrição: descreve de maneira simplificada um sistema do mundo real;
    \item Esclarecimento: aponta os principais mecanismos de um fenômeno;
    \item Experimentação: pode ser executado repetidamente, onde é possível variar os parâmetros para observar seus efeitos no comportamento do sistema como um todo e em sua saída;
    \item Prever fontes de analogia: é possível encontrar similaridades com outras simplificações, mesmo que modelem fenômenos aparentemente muito diferentes;
    \item  Comunicação/educação: pode ser usado para permitir que as pessoas explorem e entendam melhor determinados fenômenos sem a necessidade de dominar todos os detalhes do modelo, facilitando o aprendizado;
    \item Prover peças centrais para o diálogo científico: também chamado de `um objeto com que pensar', quando comparado com modelos textuais, elimina ambiguidades de interpretação, possibilitando discutir quais mecanismos são importantes para gerar determinado comportamento através da remoção ou adição de mecanismos de maneira mais objetiva;
    \item Experimento mental: não precisa necessariamente representar um fenômeno do mundo real, podendo restringir-se a um experimento mental;
    \item Predição: a partir da descrição de padrões passados, espera-se que seja possível explicar padrões futuros que possam surgir. 
\end{itemize}

Esta abordagem de construção de modelos generativos, utilizando a modelagem baseada em agentes, é também consistente com as exigências propostas por Achiles e discutidas por Rickles dentro da perspectiva da econofísica sobre quais elementos devem constituir um modelo teórico na ciência \cite{RICKLES2007948}. Para Achiles, um modelo teórico na ciência deve conter, então, cinco características:

\begin{enumerate}
    \item É um conjunto de suposições sobre algum objeto ou sistema.
    \item Essas suposições atribuem uma estrutura interna, composição ou mecanismo que se manifestam em outras propriedades exibidas pelo objeto ou sistema.
    \item Essas suposições são tratadas como uma aproximação simplificada útil para determinados fins.
    \item O modelo é proposto dentro de um arcabouço constituído por teorias mais fundamentais.
    \item O modelo pode apresentar uma analogia entre o objeto ou sistema descrito e algum outro objeto ou sistema.
\end{enumerate}

Este tipo de modelagem parte de um conjunto de suposições sobre o comportamento dos agentes (1), as quais são, por definição, aproximações simplificadas e úteis para fins de modelagem (3). A partir dessas suposições, considera-se que as propriedades emergentes do sistema resultam diretamente da estrutura interna atribuída a esse sistema (2). Estes modelos também costumam ser propostos dentro de um arcabouço teórico mais amplo, tipicamente o da mecânica estatística (4), e operam como uma analogia entre diferentes sistemas (5).

Uma vez que a abordagem de ABM é a que os físicos e os cientistas, em geral, estão mais habituados a utilizar para modelar um sistema, resulta, portanto, natural fazermos uma rápida comparação entre ambas as estratégias de modelagem: EBM e ABM. Uma diferença importante é que a ABM pode modelar populações heterogêneas de agentes com facilidade, enquanto a EBM geralmente parte da hipótese de homogeneidade (\textit{mean field}). Em muitas áreas, a heterogeneidade desempenha um papel central. Da mesma forma que a ABM pode ser utilizada para modelar interações complexas entre agentes, também pode ser utilizada para modelar interações complexas entre agentes e o ambiente, onde o ambiente é representado por agentes estacionários. Este é um caso que pode ser encontrado, por exemplo, quando utilizamos ferramentas de modelagem baseada em agentes para modelar autômatos celulares. Assim, além de a ABM ser indicada para situações em que há heterogeneidade entre os agentes, também é recomendada quando há um ambiente heterogêneo, sugerindo que, quando a heterogeneidade é uma característica importante do sistema, a ABM tende a ser a melhor escolha de modelagem.

Outra vantagem é que, enquanto para modelar com equações precisamos ter um bom conhecimento fenomenológico dos mecanismos internos do sistema de interesse, na ABM precisamos apenas das regras a nível de indivíduos, o que geralmente é mais direto de se obter. Um grande risco no processo de modelagem como um todo é modelarmos os efeitos e não as causas % \apud{phdthesis}{maturana}.
 Pela natureza da abordagem mais direta a nível individual da ABM, partimos exatamente da ideia de modelar essas interações para que o fenômeno emergente de nosso interesse seja, então, o resultado do comportamento do nosso sistema, permitindo que, assim, naturalmente, foquemos na causa para entendermos o efeito. 

Os resultados da ABM também são mais detalhados, uma vez que as saídas não são apenas a nível do sistema, mas também a nível individual. Outra diferença óbvia é que a ABM representa resultados com indivíduos de forma naturalmente discreta e frequentemente com flutuações (quando se incluem probabilidades, o que é uma prática comum não apenas quando estamos  tentando obter leis estatísticas, mas uma prática comum no uso de ABM de forma geral), enquanto a EBM fornece uma representação contínua e sem ruído do mesmo sistema.

A modelagem baseada em agentes também possui uma concepção de tempo mais rica quando comparada a outras técnicas de modelagem, pois permite observar os indivíduos interagindo ao longo do tempo. Porém, como qualquer modelo, há situações em que o custo de construir uma ABM excede seus benefícios. De maneira geral, se o problema envolve um Grande número de agentes homogêneos, será possível obter uma solução mais acurada e em menor tempo usando uma solução como campo médio. Logo, podemos considerar que este tipo de sistema é melhor modelado por esta categoria de modelos. 

Quanto ao custo computacional, há uma relação de custo e benefício semelhante à das outras formas de modelagem: resultados mais detalhados são obtidos com modelos mais complexos que exigem maior poder de processamento e armazenamento. De maneira geral, a ABM é computacionalmente mais pesada do que a EBM, uma vez que esta última é relativamente simples, pois frequentemente consiste em um pequeno conjunto de equações que podem ser resolvidas numericamente ou até mesmo por meio da análise do sistema em torno dos pontos de equilíbrio.

Por fim, vale destacar que a modelagem baseada em agentes permite manter um histórico de interação entre os agentes, e esse histórico pode servir para que o agente altere seus comportamentos e estratégias com base em eventos passados. Isso torna esta ferramenta muito útil quando se modelam interações complexas de agentes adaptativos. Devido a forma como os ABM são construídos, novas regras de interação podem ser adicionadas, removidas ou alteradas com relativa facilidade.

Como pode-se perceber, a modelagem baseada em agentes é mais um conceito para a construção de um modelo do que uma `receita'. E a escolha entre ABM e EBM depende do que e como queremos modelar; não há um método intrinsecamente superior ao outro. %PorémDentro do contexto que estamos discutindo, de sistemas complexos e leis estatísticas, estamos interessados em uma modelagem a nível microscópico que consiga ser capaz de gerar a lei estatística como uma propriedade emergente do sistema, desta forma nos é mais adequado utilizarmos ABM.

De maneira geral, pode-se dizer que a modelagem baseada em agentes é uma modelagem que utiliza agentes, que, por sua vez, são um tipo de inteligência artificial (IA); assim, é útil ter um pouco de familiaridade com esses conceitos. Inteligência artificial é o campo que visa construir sistemas que funcionarão de forma autônoma em ambientes complexos, podemos definir inteligência artificial como o estudo de sistemas que recebem entradas do ambiente e executam ações para retornar alguma saída de forma autônoma\cite{phdthesis}. %Entre teorias existentes que são utilizadas dentro de inteligência artificial, podemos mencionar especialmente o aprendizado de máquina.

Agentes autônomos %\apud{phdthesis}{maes} 
são uma abordagem no estudo da IA, fortemente inspirada na biologia. Um agente nada mais é do que um sistema que, por meio de sensores, consegue perceber o ambiente e agir sobre ele, sendo autônomo, uma vez que suas operações são realizdas sem a interferência do usuário. O sistema que forma o agente, de maneira geral, pode ser constituído de \textit{software} e \textit{hardware} permitindo, que o sistema manipule um ambiente real.

O conceito de sistema multiagente (\textit{multiagent sytem} - MAS) diz respeito a uma aplicação de agentes, onde há vários agentes, do mesmo ou de diferentes tipos, que trabalham de forma colaborativa para completar uma tarefa ou resolver um problema. Uma característica importante é que eles compartilham o mesmo ambiente, sendo necessária, portanto, uma estrutura (\textit{framework}) que permita a comunicação e a coordenação entre os agentes.

Uma simulação construída  por meio da ABM  (também chamada de \textit{agent based modelling and simulation} - ABMS) é um sistema multiagente em um ambiente simulado com um tempo virtual \cite{UFRGS}.  Dessa forma, dentro da nossa modelagem, o sistema do agente é constituído apenas de \textit{software}, sendo apenas um programa, um agente de existência virtual que percebe e opera sobre um ambiente igualmente virtual. Podemos destacar ainda que, apesar de se sobreporem muitas vezes, MAS e ABMS são tradicionalmente discutidas na literatura sob diferentes enfoques \cite{Rocha17}.  

O objetivo da ABM  costuma ser entender alguma propriedade emergente que surge do comportamento coletivo dos agentes, sendo mais usado em ciências, enquanto a MAS tem o objetivo de resolver algum problema em particular, sendo mais comumente utilizada nas engenharias. Devido a enorme variedade de aplicações, na literatura ainda podemos encontrar outros nomes para modelos construídos através da ABM, dependendo muitas vezes do enfoque, como, por exemplo, \textit{agent-based simulation} e \textit{agent-based systems} que possuem a mesma abreviatura (ABS), \textit{agent-based social simulation} (ABSS) que ainda é sub-dividido em três categorias - \textit{social aspects of agent systems }(SAAS), \textit{multi-agent based simulation} (MABS) e \textit{social simulation} (SocSim)-  além do próprio MAS, que aparece também como sinônimo em parte da literatura\cite{Rocha17}.

Antes de concluirmos esta sessão, vale destacar a importância do conhecimento de alguns princípios  básicos para uma boa modelagem. É comum a necessidade de um conhecimento matemático, pois, para uma modelagem computacional, precisamos abstrair matematicamente o sistema que queremos modelar.  Porém, outros conhecimentos específicos serão necessários e irão depender do sistema que estamos modelando. Uma vez que o cientista responsável pela modelagem precisa trabalhar com observações da natureza que fogem do seu domínio principal, passa a ser obrigatório apoiar-se em modelos já consolidados, consultar a literatura e/ou trabalhar de forma colaborativa com especialistas de outras áreas.

Uma vez que construímos um modelo satisfatório, um dos papeis mais importantes da modelagem é a simulação. Pois é através dela que podemos determinar quão próximo da realidade o modelo está. A partir do momento em que temos um grau de fidelidade aceitável no aspecto que o modelo se propõe a modelar, podemos criar situações que são difíceis de observar e estudar no mundo, cumprindo então o objetivo de que o modelo possa ser utilizado como ferramenta de tomada de decisão \cite{UFRGS}. 

\subsubsection*{Arquitetura Social do Capitalismo}


O modelo implementado nesta tese parte do modelo SA original proposto por Ian Wright~\cite{WRIGHT2005589}. No entanto, escolhemos uma notação e terminologia ligeiramente diferentes para nos alinharmos com as comumente usadas em econofísica e modelagem baseada em agentes. A sociedade é composta por $N$ agentes (rotulados como $i=1,\dots,N$ e que não são necessariamente indivíduos, mas podem representar outras entidades econômicas), em que cada agente $i$ tem uma quantidade inteira positiva $w_i(t)$ que representa a riqueza, a qual varia ao longo do tempo. % Interpretação 1: \footnote{Wright e Yakovenko enfatizam que o que é modelado é o dinheiro, e não a riqueza.} % Interpretação 2: Aqui, riqueza é definida como uma quantia de dinheiro, ativos ou bens que podem ser eventualmente transformados em dinheiro, pertencentes a cada pessoa de uma população. 
Embora essa quantidade não precise ser um número inteiro, seguimos a proposta original, pois essa distinção não é crítica para a nossa análise.

O tamanho da população, $N$, e a riqueza total do sistema, $W= \sum_{i=1}^N w_i(t)$, são conservados, portanto a riqueza per capita $\overline{w}=W/N$ é constante.  Os agentes podem pertencer a uma das três classes: empregados (classe trabalhadora), empregadores (classe capitalista) ou desempregados.  Cada agente é caracterizado por um índice $e_i\neq i$, que identifica o seu empregador: $e_i = j$ se o agente $j$ for o empregador do agente $i$, e $e_i = 0$ se o agente $i$ estiver desempregado ou for empregador\footnote{Em termos de implementação, ainda podemos aproveitar para definir $e_i<0$  para identificar facilmente quando o agente $i$ é empregador.}. Portanto, em qualquer momento $t$,  o estado de toda a economia é definido pelo conjunto de pares $S(t) = \left\{ \left( w_i(t), e_i(t) \right) :1\leq i\leq N \right\}$.
%
Uma empresa é composta por um conjunto de funcionários e um empregador, que é o único proprietário da empresa. O modelo tem quatro parâmetros originais: tamanho do sistema $N$, riqueza per capita $\overline{w}$ e os salários mínimo ($p_a$) e máximo ($p_b$) que os funcionários podem receber. Salvo indicação em contrário, todas as seleções aleatórias (agentes, alíquotas de riqueza, salários, etc.) são uniformes.

Além dos agentes, o modelo inclui um mercado $V$, representando bens e serviços produzidos pelas empresas. $V$ cresce com os gastos dos consumidores e diminui com as receitas das empresas. As empresas competem por uma parte desse fundo comum em vez de receber pagamentos diretos. 

A seguir, resumimos as regras originais, enfatizando as modificações quando existem. No estado inicial, todos os agentes têm a mesma riqueza ($w_i(0)=\overline{w}$) e o mesmo estado ($e_i(0)=0$), e o seguinte conjunto de cinco regras é aplicado, na ordem dada, após a seleção aleatória de um agente $i$.

\begin{itemize}

\item   \textbf{Contratação} -- Se $i$ estiver desempregado ($e_i=0$), um potencial empregador $j$ é selecionado do conjunto $H$  de todos os agentes que não são empregados ($e_j \neq 0$), com  probabilidade  $P (j)= w_{j}/\sum_{n\in H}w_{n}$.   Se $ w_j \ge \overline{p}$\footnote{$\overline{p}=(p_a+p_b)/2$ é o salário médio.}, o agente $i$ é contratado por $j$ (portanto, $e_i=j$).    Se $j$ estava anteriormente desempregado, $j$ torna-se um empregador ($e_j=-1$).

\item    \textbf{Despesas} (em bens e serviços) -- Um agente aleatório $j$ é selecionado e um montante $w$, representando as despesas do agente, é subtraído da sua riqueza ($w_j \to w_j-w$) e adicionado ao valor de mercado ($V \to V+w$). Modificamos esta regra, introduzindo dois novos parâmetros,  $\Omega_E$ associado à frequência de visitas ao mercado e $\Phi_E$ representando a fração da riqueza gasta no mercado. Também relaxamos a condição inicial que proibia a possibilidade de que $i=j$.
%
Se $\Omega_E\leq1$, isso denota a probabilidade de que essa regra seja aplicada.  
Caso contrário, $\Omega_E$ indica o número de vezes que a regra é aplicada. 
Para cada aplicação, um agente aleatório $j$ é selecionado e um valor de despesa aleatório $w$ é extraído do intervalo $\left[0,\Phi_E w_j\right]$. 
 
\item \textbf{Receita de mercado} (proveniente da venda de bens e serviços) -- 
Um agente aleatório $j$ é selecionado e um montante $w$ é  transferido do valor de mercado ($V \to V - w$) para o agente $j$.  Se $j$ for um empregador, o montante é creditado a $j$ ($w_j \to w_j + w$); se $j$ for um empregado, é creditado ao empregador $e_j$  ($w_{e_j} \to w_{e_j} + w$).  Em todos os casos,  $w$ representa a receita para o proprietário da empresa.

Modificamos esta regra, introduzindo dois novos parâmetros, nomeadamente,  $\Omega_M$, associado à frequência das visitas ao mercado, e $\Phi_M$, representando a fração da riqueza retirada do valor de mercado. Também relaxamos a condição de que necessariamente $i=j$.
%
Se $\Omega_M\leq1$, isso denota a probabilidade de que essa regra seja aplicada.
Caso contrário, $\Omega_M$ especifica o número de vezes que a regra é aplicada.  Para cada aplicação, um novo agente aleatório $j$ é selecionado; se $j$  não estiver desempregado ($e_i\neq 0$), um valor aleatório é extraído do intervalo $w \in [0, \Phi_M V]$.

\item \textbf{Demissão} -- Se o agente $i$ atua como empregador ($e_i=-1$), a sua restrição financeira determina quantos funcionários devem ser demitidos. Seja $n_i$ o número de agentes empregados por $i$ (ou seja, aqueles com $e_j=i$). Como cada trabalhador recebe em média um salário $\overline{p}$,  e a riqueza atual  $w_i$ do empregador permite em média sustentar uma folha salarial de  $w_i/p$ funcionários, então se $n_i > w_i/p$, o número excedente de trabalhadores

\begin{equation}
n = \max\!\left(n_i - \frac{w_i}{\overline{p}},\,0\right)
\end{equation}

são demitidos. Os $n$ agentes a serem demitidos são selecionados aleatoriamente da lista de funcionários de $i$. Se todos os funcionários forem demitidos, o agente $i$ fica desempregado ($e_i=0$).

\item  \textbf{Pagamento de salários} -- 
 Se o agente $i$ atua como empregador ($e_i=-1$), para cada agente $j$ empregado por $i$, ocorre a seguinte transferência  
\begin{center}  
  $w_i  \to w_i-p$,  \\
  $w_j  \to w_j+p$,
\end{center}
onde $p$ é um montante discreto selecionado aleatoriamente do intervalo $[p_a,p_b]$, mas se $p > w_i$, então um novo salário $p$ é selecionado do intervalo $[0,w_i]$.
\end{itemize}
  

%%%%
A atualização é assíncrona, o que significa que o estado do sistema é atualizado após a aplicação de cada regra individual. Uma sequência de $N$ aplicações de regras (um passo de Monte Carlo) corresponde, arbitrariamente, a um mês, dando a cada agente, em média, a oportunidade de ser selecionado uma vez por mês. 

Também consideramos uma variante alternativa do modelo obtida pela modificação das regras relacionadas ao emprego, inspirada também em outro trabalho de Wright\cite{WRIGHT2010}. Como estas regras vão além da inserção de um parâmetro, optamos por escrevê-las separadamente.

%%%
Propomos uma modificação na regra de  \textit{Contratação}, na qual os salários não são mais fixos, mas emergem da dinâmica do sistema. Além disso, os agentes podem procurar empregos alternativos que ofereçam salários mais altos, mesmo quando já estão empregados.
%
Para ativar esta funcionalidade, introduzimos duas variáveis específicas do agente e mais um parâmetro do sistema. As variáveis são $\mu_i$ que registra o último salário recebido pelo agente $i$, e $\eta_i$ que representa o salário esperado pelo agente $i$. O parâmetro $\alpha$ controla a pressão que rege os ajustes salariais.

Essas extensões exigem também uma reformulação das regras de  \textit{Demissão} e Pagamento de salários, além da regra de \textit{Contratação}. Estas regras estão diretamente relacionadas à dinâmica salarial. As novas regras são então:

\begin{itemize}
\item \textbf{Contratação} -- Se $i$ não é um empregador ($e_i\geq0$):
\begin{enumerate}
    \item Um potencial empregador $j$ é selecionado do conjunto $H$  de todos os agentes, exceto funcionários ($e_j\neq0$), com probabilidade  $P (j)= w_{j}/\sum_{n\in H}w_{n}$.

     \item Uma oferta salarial $w$ é retirada uniformemente do intervalo $\left[\eta_{i},\left(1+\alpha\right)\eta_{i}\right]$,    onde $\eta_i$ representa a expectativa salarial do agente e o parâmetro $\alpha$ controla, portanto, a intensidade do aumento de expectativa salarial pelo agente.  Se $w$ exceder a riqueza do empregador selecionado ($w>w_j$), é truncado para $w=w_j$.

    \item Se a oferta salarial final exceder a expectativa do agente ($w > \eta_i$), o agente $i$ é contratado pelo empregador $j$ ($e_i = j$), e a expectativa é atualizada para $\eta_i = w$. Se o agente $i$ não for contratado e estava originalmente desempregado, uma nova expectativa é obtida uniformemente a partir do intervalo $\left[0, \eta_i\right]$.
    %
  
\end{enumerate}

\item \textbf{Pagamento de salários e demissões} -- Combinado agora em uma única regra, se o agente $i$ é um empregador ($e_i < 0$), definimos o conjunto $E$ de funcionários de $i$ como aqueles agentes $j$ para os quais $e_j = i$. Para cada funcionário $j \in E$:


\begin{enumerate}
    \item Se o empregador $i$ tiver fundos suficientes, um montante igual à expectativa do empregado $\eta_j$ é transferido do empregador ($w_i \rightarrow  w_i- \eta_j$) para o empregado ($w_j \rightarrow w_j +\eta_j$), e o último salário recebido pelo empregado é atualizado para $\mu_j = \eta_j$.
    
    \item Se o empregador $i$ não tiver fundos suficientes, o empregado $j$ é despedido e a sua expectativa salarial é redefinida para o último salário recebido, $\eta_j = \mu_j$. Ou seja, se o agente for despedido antes de receber o salário negociado, a expectativa não é concretizada. 

    \item Se, após processar todo o conjunto $E$, todos os funcionários tiverem sido demitidos, o empregador $i$ ficará desempregado ($e_i = 0$).

\end{enumerate}
\end{itemize} 

Os dados são coletados anualmente ao final do ano. O PIB é definido como a receita agregada de todas as empresas ao longo de um ano (12 meses) através da aplicação da regra de \textit{Receita de mercado}. O salário anual de cada agente inclui toda a renda recebida ao longo do ano, originada pela aplicação da regra de \textit{Pagamento de salários}. De modo análogo, a receita dos capitalistas corresponde à renda obtida ao longo do ano através da aplicação da regra \textit{Receita de mercado}, lembrando que, neste modelo, não há distinção entre a riqueza da empresa e a do capitalista dono da empresa. A renda total anual do agente é obtida, então, somando as receitas que obteve como capitalista e os salários que recebeu como trabalhador.


\subsection{Econofísica}

Uma vez que compreendemos a abordagem de investigação de sistemas complexos utilizando leis estatísticas e modelos microscópicos,e também entendemos o conceito de modelagem baseada em agentes que será utilizado para construir os modelos generativos, podemos buscar definir um pouco melhor o que é a econofísica em si. Há diversos trabalhos clássicos que podem ser consultados para aprofundar os estudos \cite{cottrell2009classical,slanina2013essentials,Yakovenko2020}\footnote{Yakovenko também publicou um estudo interessantíssimo acerca da economia monetária sob uma perspectiva econofísica \cite{Yakovenko2016} que merece ser mencionado.}, além de um relatório geral de Física Social \cite{JUSUP20221}. Para mais recomendações, é possível também consultar o trabalho de Rickles no qual 11 outros livros são avaliados\cite{RICKLES2007948}, assim como o artigo resenha do economista Fernando Nogueira da Costa\cite{Fernando}.

Mas devemos, pelo menos, esboçar uma ideia geral do que se trata a econofísica. Esta ideia de que a física poderia ser utilizada para predizer como grupos humanos se comportam é uma possibilidade que é popularmente mais conhecida na ficção científica. Podemos citar como exemplo a trilogia do escritor russo-americano Isaac Asimov chamada Fundação. Nesta obra, o protagonista apresenta uma ciência chamada ``psico-história". Esta ciência seria capaz de prever computacionalmente o comportamento da sociedade, pois, apesar de o comportamento individual ser imprevisível, o comportamento de multidões poderia ser tratado estatisticamente \cite{parthum}.

Como discutido anteriormente, esta ideia não ficou restrita a ficção. Algo que, nos dias atuais, pode desafiar o senso comum consiste no fato histórico de que os métodos de probabilidade e estatística floresceram primeiro entre os cientistas sociais que investigavam regularidades quantitativas no interior do funcionamento da sociedade; será somente em um segundo momento que a própria física se move do determinismo da física Newtoniana em direção à descrição probabilística de gases ideais.

%Augusto Comte no começo do século 19 foi o  responsável por cunhar o termo ``física social'' como uma referência explícita ao sucesso da mecânica Newtoniana. Apesar dele próprio ter abandonado o termo, ele foi recuperado mais tarde por seus sucessores, particularmente Adolphe Quételet, astrônomo por formação, que publicou, em 1835 um ensaio sobre `física social'.   Outro intelectual influente para o desenvolvimento da área foi o historiador Henry Buckle, responsável por apresentar uma abordagem estatística da história  da civilização.  

%Esses conceitos desenvolvidos inicialmente fora da física e inspirado pela física, terminam por influenciar diretamente o próprio desenvolvimento da física. 

Nos dias atuais, os físicos recorrem a analogia entre moléculas e humanos para investigar questões sociais; por outro lado, Boltzmann tinha em mente uma analogia inversa quando estava investigando os gases. Boltzmann  buscou inspiração na sociedade para analisar um fenômeno tipicamente físico \cite{parthum}. Em um trecho do seu artigo de 1872, ele escreve:   ``As moléculas são como muitos indivíduos, tendo os mais diversos estados de movimento, e as propriedades dos gases permanecem inalteradas só porque o número destas moléculas, que em média têm um dado estado de movimento, é constante\cite{boltzmann1872_translation}'' (em tradução livre).

Uma desconfiança genuína que pode surgir diz respeito a simplicidade da modelagem de um fenômeno tão complexo como a sociedade humana. Porém, devemos lembrar que, em qualquer processo de modelagem, sempre há um grande número de detalhes negligenciados. Todo modelo é sempre  uma simplificação da realidade; o que se busca ao modelar um sistema social, assim como em outras áreas mais tradicionais da ciência, é incorporar não todos os detalhes do  fenômeno real, mas os ingredientes essenciais que nos permitirão descobrir e investigar os mecanismos fundamentais do fenômeno que estamos observando. 

Um modelo com detalhes demais pode até mesmo ser contraproducente para este objetivo, pois o excesso de detalhes pode mascarar esses ingredientes essenciais. É também uma característica intrínseca da própria física estatística a existência de uma certa independência em relação a todos os detalhes dos elementos que constituem o sistema a ser investigado; em muitos casos, basta que conheçamos algumas poucas características fundamentais do sistema  para entender as estruturas emergentes desse sistema. Como consequência disso, temos que um coletivo pode apresentar certa previsibilidade, apesar de o comportamento individual dos agentes que formam esse coletivo ser imprevisível.

Outro intelectual influente para o desenvolvimento da área foi o historiador Henry Buckle (1821-1862), responsável por apresentar uma abordagem estatística da história  da civilização. Georg Graf (1781-1851) contemplou  a possibilidade de aplicar os princípios da mecânica clássica especificamente na economia; seu conhecimento da física newtoniana era profundo e ele contribuiu para diversas áreas da ciência. Um texto seu, datado de 1815, já apresentava o que poderia ser um capítulo de um livro de economia e se referia explicitamente a mecânica clássica; essa forma de pensar era revolucionária para o seu tempo. 

Na segunda metade do século XIX, uma tentativa similar foi feita por Vilfredo Pareto (1848-1923). Pareto também se voltou para a economia, munido de um bom conhecimento de mecânica clássica newtoniana e de um bom entusiasmo para descrever movimentos sociais da mesma forma que descrevia o movimento dos planetas, comparando inclusive as leis de Kepler às leis econômicas ainda a serem descobertas. A lei de Pareto de distribuição da riqueza é agora uma peça padrão do conhecimento econômico; porém, sua redução das leis econômicas a algo similar às leis de Newton permaneceu como um mero sonho até os dias de hoje.

A física mudou completamente com a descoberta da teoria da relatividade e da mecânica quântica no início do século XX. A questão que naturalmente surgiu foi se a velha mecânica newtoniana era inadequada para explicar sistemas complexos como a sociedade e a economia; o mesmo se aplica à nova física de Einstein e Schrodinger? Uma nova física abre novos caminhos para fazer essa aproximação de maneira científica? Ettore Majorana (1906-1938), físico teórico, tornou-se popular ao realizar uma discussão sobre essas questões. Benoit B. Mandelbrot (1924-2010) deu sua contribuição histórica para a interação entre a física e a economia, propondo a introdução do conceito de fractais. Como podemos notar, sistemas complexos parecem fadados a perseguir uma rota diferente da física tradicional.

Um marco decisivo na história do desenvolvimento da econofísica  foi uma conferência programada para ocorrer em Moscou entre 1 e 5 de julho de 1974, na qual a organização do comitê incluía pessoas como Kenneth Arrow (1921-2017) e Hans Bethe (1906-2005), um laureado em economia e outro em física, respectivamente. Este evento, apesar de ter sido cancelado, desempenhou um papel significativo na transferência das ideias e da linguagem da física para outros ramos do conhecimento humano. Seus  manuscritos sobreviveram e alguns foram publicados posteriormente, influenciando diretamente o desenvolvimento futuro da área\cite{slanina2013essentials}.

Após este evento, outras tentativas de misturar física e outras disciplinas começaram a florescer ainda mais fortemente. Entre os anos 70 e 80, a ciência da complexidade começou a se tornar imensamente popular, com o Instituto Santa Fé tendo um papel ativo neste processo. Foi ali  que a teoria de sistemas complexos começou a ser sistematicamente aplicada na economia. P. W. Anderson (1923-2020), outro laureado em física, começou a se envolver profundamente e a promover esta linha de pesquisa. Em 1991, a Physics A publicou o que é considerado agora como o primeiro artigo que pode ser verdadeiramente atribuído ao recém nascido campo da econofísica.

Mas, afinal, o que é exatamente a econofísica? Esta é uma discussão que não possui consenso sobre a resposta. Para Slanina, não é uma simples mistura entre ciências naturais e economia, também não é uma ciência interdisciplinar entre economia e física\cite{slanina2013essentials}. Ele opta, então, por chamá-la de transdisciplinar, pois  não quer encontrar um lugar entre as disciplinas, mas sim encontrar princípios que são verdadeiros em ambas as disciplinas. Isto é, uma busca por princípios comuns, ferramentas comuns e consequências comuns. E após encontrarmos princípios comuns, devemos, então, checar se esses princípios são úteis.

Esta questão nos remete diretamente a outra: de que forma exatamente os econofísicos empregam seus esforços? Slanina nos diz que os métodos provenientes da física são usados sistematicamente para desenvolver análises empíricas e na construção de modelos, usando técnicas matemáticas avançadas para aplicar no mundo real; isto é , a econofísica não é sobre pegar um fragmento da física teórica e reinterpretá-lo na língua da economia; ela tenta ser abrangente e coerente.

Já o relatório de física social de 2022 adota uma definição mais ampla: física social é uma coleção de tópicos de pesquisa que visam resolver problemas sociais onde cientistas com treinamento formal em física tem contribuído e continuam a contribuir significativamente\cite{JUSUP20221}. Esta é uma definição pragmática baseada na ideia de que o que os físicos realmente fazem é o que define o que é física. O enorme sucesso da física no último século levou cientistas de outras disciplinas a tentarem formular métodos quantitativos similares aos da física; isto é, de maneira geral, houve uma proliferação de modelos matemáticos e físicos em diversas disciplinas. E então, se vamos utilizar métodos análogos aos da física, quem melhor do que os próprios físicos para contribuir? Dessa forma, opta-se por chamar de física social esta coleção de tópicos que tem como objeto de estudo problemas sociais e que conta com os físicos como os profissionais empenhados em tentar investigar.  

Embora a física social não seja algo novo, essa tendência ressurgiu com mais intensidade nos tempos atuais. Isso também surge como resultado da diminuição no ritmo do progresso científico na física tradicional, ao menos em comparação com o século anterior. Esse fenômeno resulta, então, por um lado, no movimento de alguns físicos que partiram em busca de uma `nova física' e, por outro lado, no de jovens cientistas que decidiram empregar suas habilidades quantitativas em outros tópicos. Estes tópicos da física social muitas vezes surgem a partir de uma reflexão sobre o que constitui e permite o modo de vida moderno, e o que perturba e ameaça esse modo de vida. A prosperidade das cidades esta de muitas formas atrelada aos mercados; a economia tem um enorme impacto no desenvolvimento e planejamento da vida urbana, e assim surge especificamente o tópico de econofísica no interior da própria física social.

Já Rickles discute a existência de diferentes definições de econofísica, classificando-as em quatro grupos diferentes \cite{rickles2023econophysics}:
\begin{itemize}
    \item Orientada para o conceito: a econofísica envolve uma espécie de mecânica estatística da economia.
    \item Orientada para a metodologia: a econofísica é o estudo empírico adequado da economia, utilizando uma abordagem que privilegia os dados.
    \item Orientada para a ontologia: a econofísica envolve a visão dos sistemas econômicos como sistemas complexos, com agentes econômicos interagindo como componentes, de modo que os sistemas econômicos seguem, de maneira mais geral, leis semelhantes às dos sistemas ''naturais``.
    \item Orientada para a heurística: a econofísica envolve encontrar analogias entre sistemas financeiros/socioeconômicos e físicos (naturais).
\end{itemize}

Os econofísicos constroem suas teorias em torno da replicação de alguns conjuntos de fatos econômicos. Esses fatos são propriedades comuns ao longo de diferentes mercados e períodos de tempo e são conhecidos pelo nome de fatos estilizados. Esses fatos estilizados, por sua vez, são obtidos por meio da abstração dos dados empíricos e constituem o núcleo central da metodologia de pesquisa adotada pela econofísica, uma vez que o objetivo da construção de modelos econofísicos é exatamente fornecer uma explicação estrutural para um ou mais fatos estilizados.  Temos, por exemplo, a desigualdade observada na distribuição de renda, que constitui um fato estilizado da economia, cuja regularidade é frequentemente descrita por uma lei estatística, sendo a lei de Pareto uma das mais utilizadas.

A econofísica, historicamente, têm se focado em questões como o fluxo de dinheiro e os mercados financeiros, devido  a existência de um grande conjunto de dados disponíveis que permitiu que a aplicação de ferramentas e métodos da física estatística se mostrasse útil nessas áreas.  Áreas como a macroeconomia  historicamente não atraíram a atenção dos econofísicos. Porém, com o crescimento da área, o aumento da disponibilidade de dados e o aumento do poder computacional, os econofísicos têm se engajado cada vez mais nesta pesquisa, expandindo o campo de atuação da econofísica das finanças para a economia. Como exemplo do crescente interesse por problemas  de economia, e não apenas finanças, podemos citar os livros seminais  de Farjoun e Machover \cite{farjoun2020laws,farjoun2022labor}, que, mesmo sem se definirem como econofísica, propuseram uma abordagem probabilística da economia inspirada na física. Também podemos citar a chegada de novos livros no mercado que propõem explicitamente a investigação da macroeconomia através de uma abordagem econofísica\cite{aoyama2017macro}.

%\subsection{Economia e econofísica}

Não é possível encerrarmos a discussão sobre econofísica sem tecermos alguns comentários sobre sua relação com a abordagem ortodoxa da economia.  Para Rickles, a construção microscópica do fato estilizado de baixo para cima, visando a compreensão da constituição desse fenômeno, cria um flagrante contraste com a abordagem de cima para baixo adotada pela economia ortodoxa\cite{rickles2023econophysics}.  A centralidade que a econofísica atribui aos fatos estilizados, modelados como leis estatísticas a partir dos dados empíricos, torna sua metodologia fortemente orientada por dados uma característica marcante.

%Dessa forma, a econofísica teria uma preocupação que tradicionalmente não se encontra na abordagem ortodoxa da economia quanto às estruturas internas que geram os fenômenos emergentes \cite{RICKLES2007948}.

Outra diferença fundamental diz respeito à racionalidade dos agentes. Enquanto a abordagem tradicional da economia visa construir modelos baseados em escolhas racionais, a econofísica se destaca por inicialmente tomar emprestado o modelo de partículas da física estatística para explicar o comportamento dos agentes. Esta forma de modelagem assume que os interesses e preferências dos agentes não são fixos, mas dependem da interação com outros agentes. Assim sendo, a econofísica coloca uma grande ênfase no ambiente social no qual o agente está inserido\cite{JUSUP20221}.

Neste sentido, Fernando Nogueira da Costa comenta que, na economia ortodoxa, a 'regra' é um reducionismo atomístico, onde a realidade deve ser explicada em termos de um agente representativo racional. Assume-se que o comportamento agregado do sistema é idêntico à soma dos efeitos de cada causa individual. Já a abordagem não ortodoxa da economia, que a trata como um sistema complexo, tem por `princípio' o reducionismo interativo; isto é,  ela não tenta descrever um agente de forma atomizada, mas sim descrever a interação desse agente com os outros. Nesta abordagem, a descrição estatística do comportamento agregado do sistema é resultado da interação entre agentes \cite{Fernando}. Fernando Nogueira da Costa insere a econofísica no conjunto das abordagens não ortodoxas da economia, agrupadas sob a denominação de economia da complexidade, caracterizadas pela compreensão da economia como um sistema complexo.


No modelo padrão para a construção de teorias econômicas, atribui-se a cada agente uma função  de utilidade que descreve suas preferências e como maximizar sua utilidade ou felicidade; tudo isso é expresso matematicamente. Economistas tradicionais, então, usam computadores para resolver essas equações e não simulações, no sentido moderno. Uma abordagem que difere da adotada pela econofísica, conforme exposto previamente.

Uma abordagem não ortodoxa, baseada em sistemas complexos e na modelagem baseada em agentes, também induz a uma forma totalmente diferente de lidar com as mudanças que ocorrem no interior de um sistema econômico. Nesta abordagem, as mudanças não surgem por causas externas, mas devido a causas internas ao próprio sistema, quando as forças do mercado agem sobre si mesmas. Isto difere novamente da abordagem tradicional da economia, que tende a atribuir causas exógenas às crises.

%Opção 1
Podemos concluir que a econofísica pode ser compreendida como um conjunto específico de abordagens não ortodoxas da economia, que concebe o sistema econômico como um sistema complexo. Sua especificidade reside, por um lado, em seu histórico de inspiração na física e, por outro, principalmente, em uma orientação fortemente baseada em dados, voltada à reconstrução e explicação dos fatos estilizados.

% Opção 2
%Podemos concluir que a econofísica pode ser compreendida como um conjunto específico de abordagens não ortodoxas no interior da economia, que concebem o sistema econômico como um sistema complexo. Sua especificidade reside, por um lado, em sua origem histórica e inspiração metodológica na física e, por outro, sobretudo, em uma orientação fortemente baseada em dados, voltada à identificação, reconstrução e explicação de fatos estilizados enquanto regularidades estatísticas empiricamente observadas.

%É valido mencionar que a divisão da economia entre ortodoxos e heterodoxos nem sempre é clara, nesse sentido \cite{lavoie2014postkeynesian}, propõe dividir a economia em dois grandes blocos: convencional (\textit{mainstream}) e dissidentes, sendo que todo dissidente é heterodoxo, mas nem todo heterodoxo deixa de ser convencional. 

Esta pesquisa adota uma metodologia baseada na concepção da economia como um sistema complexo, utilizando modelagem computacional baseada em agentes. Modelos microscópicos são construídos com o objetivo de gerar, como fenômenos emergentes, regularidades macroscópicas, conhecidas como fatos estilizados, observadas empiricamente e expressas na forma de leis estatísticas extraídas de dados reais. A validação ocorre por meio da comparação estatística entre os resultados das simulações e os dados empíricos.

%
\subsection{Lei de Pareto como lei estatística}
Além da construção do modelo generativo que discutimos com mais detalhes previamente, outra atividade necessária na investigação de sistemas complexos adotando a abordagem proposta consiste na construção de ao menos uma lei estatística macroscópica que deve emergir deste modelo generativo a ser construído. Nesta pesquisa, não nos ocupamos desta parte da investigação; porém, caso o leitor esteja interessado, o livro elaborado por Altmann dedica um capítulo a esta questão, isto é, Como leis estatísticas podem (e devem) ser obtidas a partir dos dados disponíveis\cite{Altmann_2024}.


%{\color{red} talvez INSERIR A DISCUSSÃO DA LIN LIN SOBRE DIFERENTES DISTRIBUIÇÕES E SER BIMODAL}

Segundo Altmann\cite{Altmann_2024}, a lei de desigualdade de Pareto é a mais influente e provavelmente o primeiro exemplo de lei estatística como uma distribuição  de lei de potência.  Ela surgiu de uma análise empírica da distribuição de renda de diferentes países, buscando identificar qual a proporção $N\left(x\right)$ da população possui uma renda maior que $x$. Os dados que foram compilados por Pareto de fontes originais são principalmente do século XIX, e o resultado de sua análise foi proposto da seguinte forma:
$$\ln N\left(x\right)=\ln A-\tilde{\gamma}\ln\left(x\right) \label{pareto}$$
Com o mesmo valor $\tilde{\gamma} \approx 1.5$ sendo observado em cenários completamente diferentes, Pareto também percebeu que essa aproximação lembra uma função de distribuição cumulativa complementar (\textit{Complementary cumulative distribution function} - CCDF). A alegação de que esta distribuição descreve a renda (e riqueza) em diferentes países e regiões é chamada de lei de Pareto.  

O quão bem os dados se ajustam a esta lei é algo de debate desde que a proposta foi feita pela primeira vez.  Além de uma lei de potências, também foram feitas propostas de que os dados seriam bem descritos por distribuições do tipo Lognormal (ou Lognormal generalizado) ou pela lei exponencial de Boltzmann-Gibbs. Pesquisas empíricas mais recentes tem indicado que a distribuição de renda pode ser separada em uma região de baixa renda e uma região de maior renda, passando de uma distribuição exponencial para uma distribuição de lei de potência. De maneira geral, pesquisas recentes examinando os dados fornecidos pela World Distribution of Income tem sugerido uma distribuição bimodal \cite{diss_mods_00007721}.

%Linlin também propõe uma distribuição multmodal.

De toda forma, mais de um século depois de ter sido proposta, a utilidade das distribuições do tipo Pareto para caracterizar distribuições de renda e riqueza continua a ser reconhecida nas análises econômicas modernas. Uma vez que estamos em posse da lei estatística, conforme discutimos, um caminho possível é procurar a explicação da origem da lei, uma tarefa que o próprio Pareto se dedicou inicialmente. Esta tarefa continua a ser desenvolvida até os dias de hoje, uma vez que nenhuma explicação foi capaz de fazer com que os cientistas considerassem este um caso encerrado. 

O outro caminho possível é analisar suas consequências. Sobre isso, um debate que tem sido realizado é como a variação ou a falta de variação do  exponente $\tilde{\gamma}$ afetaria o bem estar e a desigualdade na população; assim como foram realizadas discussões econômicas sobre como melhorá-lo, assumindo a validade da lei. Evidentemente, todas essas discussões estão imersas em controvérsias.

%Joseph Persky, Retrospectives: Pareto’s Law, Journal of Economic Perspectives 6 (1992), no. 2, 181–192. 

  
\subsubsection*{Definição matemática}
A lei de Pareto nos dá a probabilidade de uma pessoa ter exatamente o valor $x$ da variável de interesse, no nosso caso, a probabilidade de um agente ter uma riqueza no valor $w$ . A função de densidade de probabilidade (\textit{Probability Density Function} - PDF)\footnote{Esta função é usada para especificar a probabilidade de uma variável aleatória estar dentro de um determinado intervalo de valores.}  conhecida como lei de Pareto, é \cite{pedroso}:

$$p\left(w\right)=\frac{\alpha\beta^{\alpha}}{w^{\alpha+1}}$$

%ebta também é riqueza mínima
Onde $\alpha >0$ é o coeficiente de Pareto (parâmetro de forma) que modela a distribuição e $\beta >0$ é a riqueza mínima (parâmetro de escala). O parâmetro que normalmente domina as discussões sobre a lei de Pareto refere-se ao índice de Pareto. Dessa forma, um objetivo comum é calcular o coeficiente de Pareto nas distribuições que estamos analisando. Para isso, iremos trabalhar com a função de distribuição cumulativa complementar.  

A função de distribuição cumulativa (\textit{Cumulative Distribution Function }- CDF) de uma variável aleatória de valor real $x$ é a função dada por:
$$F_{X}\left(x\right)=P\left(X\leq x\right)$$
Ou seja, temos a probabilidade de uma variável aleatória $X$ assumir um valor menor ou igual a $x$. A CDF de uma variável contínua $X$  pode ser expressa em termos de uma função densidade $f_x$ de probabilidade dada por:
$$F_{X}\left(x\right)=\int_{-\infty}^{x}f_{X}\left(t\right)dt$$

A CCDF pode ser definida, então, como:
$$\overline{F}_{X}\left(x\right)=P\left(>x\right)=1-F_{X}\left(x\right)$$
Isto é, a probabilidade de uma variável aleatória $X$ assumir um valor maior que  $x$. No nosso caso,  estamos interessados na probabilidade de que, ao selecionarmos um agente aleatoriamente, ele tenha uma riqueza superior a $w$, fazendo uma pequena mudança de notação, temos:

$$\overline{F}_X\left(w\right)=P\left(X>w\right)=1-F_{X}\left(w\right)$$

Podemos calcular ,então, a CDF de Pareto da seguinte forma:


$$F_{X}\left(w\right)=Pr\left(X\leq w\right)=\int_{-\infty}^{w}p\left(w'\right)dw'$$

Resolvendo a integral:
\begin{equation}
    \begin{split}
        F_{X}\left(w\right)	& =\int_{\beta}^{w}\frac{\alpha\beta^{\alpha}}{x'^{\alpha+1}}dw' \\    & =\alpha\beta^{\alpha}\left(\frac{w'^{-\left(\alpha+1\right)+1}}{-\left(\alpha+1\right)+1}\right)|_{\beta}^{w}\\
    	&=\alpha\beta^{\alpha}\left(\frac{w'^{-\alpha}}{-\alpha}\right)|_{\beta}^{w}\\
    	&=\alpha\beta^{\alpha}\left(\frac{w^{-\alpha}}{-\alpha}-\frac{\beta^{-\alpha}}{-\alpha}\right)\\
    	&=\left(\frac{\beta}{\beta}\right)^{\alpha}-\left(\frac{\beta}{w}\right)^{\alpha}\\
        &=1-\left(\frac{\beta}{w}\right)^{\alpha}
    \end{split}
\end{equation}
Então, o CCDF da lei de Pareto será:

$$\overline{F}_X\left(w\right)=1-\left(1-\left(\frac{\beta}{w}\right)^{\alpha}\right)=\left(\frac{\beta}{w}\right)^{\alpha}=aw^{-\alpha}$$

Onde $a=\beta^{\alpha}$.  Se aplicarmos o logaritmo natural, temos: 

$$\ln\overline{F}_{X}\left(w\right)	=\ln\left(aw^{-\alpha}\right)=\ln\left(a\right)+\ln\left(w^{-\alpha}\right)=\ln\left(a\right)-\alpha\ln\left(w\right)$$
Onde retomamos a equação \ref{pareto}, identificando $\overline{F}(w) =N(x)$  e  $\alpha = \tilde{\gamma}$. 

Porém, estamos interessados em obter o índice de Pareto de uma distribuição de valores. Isso é obtido mais facilmente após realizarmos uma linearização. Para linearizar o CCDF da lei de Pareto, começamos aplicando o logaritmo:

\begin{equation}
    \begin{split}
         \log_{10}\left(\overline{F}_X\left(w\right)\right) & =      \log_{10}\left(aw^{-\alpha} \right)\\      
        	&=\log_{10}\left(a\right)+\log_{10}\left(w^{-\alpha}\right) \\
    \end{split}
\end{equation}
Como $a$ é um valor constante, Nada nos impede de escrever $a=10^{b}$ . Escrevendo a linearização do CCDF de Pareto como  $y=\log_{10}\left(\overline{F}_X\left(w\right)\right)$, temos então:

\begin{equation}
    \begin{split}
        y	     &=\log_{10}\left(10^{b}\right)+\log_{10}\left(w^{-\alpha}\right) \\
        &=b\log_{10}\left(10\right)-\alpha\log_{10}\left(w\right)
    \end{split}
\end{equation}
Fazendo  uma substituição de variável do tipo $x=\log_{10}\left( w \right)$ , ficamos então com a seguinte equação linear:

$$y=b-\alpha x$$
O índice de Pareto nada mais é do que a inclinação desta reta. Além disso, podemos escrever a constante $b$  em termos das constantes  originais da lei de Pareto da seguinte forma:

$$b=\log_{10}\left(a\right)=\log_{10}\left(\beta^{\alpha}\right)$$
Ou ainda, de forma inversa, podemos obter $\beta$  a partir da constante $b$ e do coeficiente $\alpha$ calculando $\beta=10^{\left(b/\alpha\right)}$.  Este é o CCDF linearizado da lei de Pareto , após aplicarmos o logaritmo na base 10, isto precisará ser retomado mais tarde. Para calcular o índice de Pareto de um conjunto de dados, precisaremos aplicar algum método matemático a esse conjunto de dados.  Aplicaremos a  regressão linear.

\subsection*{Regressão Linear}

Podemos aplicar o método da regressão linear  sobre um grande conjunto de dados para obter a reta que melhor descreve esse conjunto de dados. Evidentemente, para que esse método seja útil, precisamos que exista uma reta que descreva os dados de forma satisfatória. 

Se tivermos uma única variável independente $y$ que esperamos que seja descrita por uma única variável dependente $x$ de forma linear (uma linha reta), então, para um conjunto de pontos de entrada $\left\{ x_{i}\right\}$, queremos encontrar um conjunto de pontos de saída $\left\{ y_{i}\right\}$ que pode ser escrito da seguinte forma:

$$y_{i}=b_{0}+b_{1}x_{i}+e_{i}$$
onde:

\begin{itemize}
    \item $b_0$  é a constante de regressão;
    \item $b_1$  é o coeficiente de regressão;
    \item $e_i$  é o erro que obtemos quando tentamos predizer  $y_{i}$  a partir de  $x_{i}$.
\end{itemize}

Quando aplicamos o modelo de regressão linear, queremos minimizar o erro. Para fazer isso, precisamos de uma medida do erro. Usaremos a soma dos quadrados dos erros:

$$E=\sum_{i=0}^{N}e_{i}^{2}$$
Considerando que temos $N$ pontos. Então, queremos encontrar $b_{0}$ e $b_{1}$ que minimizem essa quantidade. Podemos reescrever isso como:

$$E=\sum_{i=0}^{N}\left(y_{i}-b_{0}-b_{1}x_{i}\right)^{2}$$

Podemos usar o gradiente para encontrar o mínimo de uma função. O mínimo deve satisfazer $\nabla f=0$. Como queremos encontrar o mínimo para ambos os termos $b_j$, temos:

$$\nabla E=\left(\frac{\partial E}{\partial b_{0}},\frac{\partial E}{\partial b_{1}}\right)$$

Calculando a derivada para $b_0$:
\begin{equation}
    \begin{split}
        \frac{\partial E}{\partial b_{0}}	& =\frac{\partial}{\partial b_{0}}\sum_{i=0}^{N}\left(y_{i}-b_{0}-b_{1}x_{i}\right)^{2} \\
    	&=\sum_{i=0}^{N}\frac{\partial}{\partial b_{0}}\left(y_{i}-b_{0}-b_{1}x_{i}\right)^{2}\\
    	&=\sum_{i=0}^{N}2\left(y_{i}-b_{0}-b_{1}x_{i}\right)\frac{\partial}{\partial b_{0}}\left(y_{i}-b_{0}-b_{1}x_{i}\right)\\
    	&=\sum_{i=0}^{N}2\left(y_{i}-b_{0}-b_{1}x_{i}\right)\left(-1\right)\\
    	&=-2\sum_{i=0}^{N}\left(y_{i}-b_{0}-b_{1}x_{i}\right)
    \end{split}
\end{equation}

Procedendo de modo similar para $b_1$:

$$\frac{\partial E}{\partial b_{1}}=-2\sum_{i=0}^{N}x_{i}\left(y_{i}-b_{0}-b_{1}x_{i}\right)$$

Para minimizar, precisamos que a seguinte condição seja satisfeita:

$$\nabla E=\left(\frac{\partial E}{\partial b_{0}},\frac{\partial E}{\partial b_{1}}\right)=\left(0,0\right)$$

Então:

\begin{equation}
    \begin{split}
        0&	=-2\sum_{i=0}^{N}\left(y_{i}-b_{0}-b_{1}x_{i}\right)\\
        0	&=-2\sum_{i=0}^{N}x_{i}\left(y_{i}-b_{0}-b_{1}x_{i}\right)\\
    \end{split}
    \label{linear 1}
\end{equation}

Para ter certeza de que se trata de um mínimo e não de um máximo, precisamos verificar se a segunda derivada é positiva:

\begin{equation}
    \begin{split}
        \frac{\partial^{2}E}{\partial b_{0}^{2}}	& =2\sum_{i=0}^{N}1=2N \\
        \frac{\partial^{2}E}{\partial b_{1}^{2}}	& =2\sum_{i=0}^{N}x_{i}x_{i}=2\sum_{i=0}^{N}x_{i}^{2}
    \end{split}
\end{equation}

Como $N>0$  e $x^2\geq0$, temos então a certeza de que é um mínimo.  Manipulando a primeira equação em \ref{linear 1}, obtemos: 

\begin{equation}
    \begin{split}
        0	& =-2\sum_{i=0}^{N}\left(y_{i}-b_{0}-b_{1}x_{i}\right)\\
        0 & =\sum_{i=0}^{N}y_{i}-\sum_{i=0}^{N}b_{0}-\sum_{i=0}^{N}b_{1}x_{i}\\
        b_{0}\sum_{i=0}^{N}1	& =\sum_{i=0}^{N}y_{i}-\sum_{i=0}^{N}b_{1}x_{i}\\
        Nb_{0}	& =\sum_{i=0}^{N}y_{i}-\sum_{i=0}^{N}b_{1}x_{i}\\
        b_{0}	& =\sum_{i=0}^{N}\frac{y_{i}}{N}-b_{1}\sum_{i=0}^{N}\frac{x_{i}}{N}\\
        b_{0}	& =\overline{y}-b_{1}\overline{x}
    \end{split}
\end{equation}

E para a segunda equação em \ref{linear 1} temos:
 \begin{equation}
     \begin{split}
        0	&=-2\sum_{i=0}^{N}x_{i}\left(y_{i}-b_{0}-b_{1}x_{i}\right)\\
        0&	=\sum_{i=0}^{N}x_{i}y_{i}-b_{0}\sum_{i=0}^{N}x_{i}-\sum_{i=0}^{N}b_{1}x_{i}^{2}\\
        b_{1}\sum_{i=0}^{N}x_{i}^{2}	&=\sum_{i=0}^{N}x_{i}y_{i}-b_{0}\sum_{i=0}^{N}x_{i}         
     \end{split}
 \end{equation}
 
Combinando os resultados, ficamos com:

\begin{equation}
    \begin{split}
        b_{1}\sum_{i=0}^{N}x_{i}^{2}	& =\sum_{i=0}^{N}x_{i}y_{i}-b_{0}\sum_{i=0}^{N}x_{i} \\
        b_{1}\sum_{i=0}^{N}x_{i}^{2}	& =\sum_{i=0}^{N}x_{i}y_{i}-\left(\sum_{i=0}^{N}\frac{y_{i}}{N}-b_{1}\sum_{i=0}^{N}\frac{x_{i}}{N}\right)\sum_{i=0}^{N}x_{i} \\
        b_{1}\sum_{i=0}^{N}x_{i}^{2}	& =\sum_{i=0}^{N}x_{i}y_{i}-\frac{1}{N}\left(\sum_{i=0}^{N}y_{i}\right)\left(\sum_{i=0}^{N}x_{i}\right)+\frac{b_{1}}{N}\left(\sum_{i=0}^{N}x_{i}\right)^{2} \\
        b_{1}\left(\sum_{i=0}^{N}x_{i}^{2}-\frac{1}{N}\left(\sum_{i=0}^{N}x_{i}\right)^{2}\right)	& =\sum_{i=0}^{N}x_{i}y_{i}-\frac{1}{N}\left(\sum_{i=0}^{N}y_{i}\right)\left(\sum_{i=0}^{N}x_{i}\right)
    \end{split}
\end{equation}
Então $b_1$ é:

$$b_{1}=\frac{\sum_{i=0}^{N}x_{i}y_{i}-\frac{1}{N}\left(\sum_{i=0}^{N}y_{i}\right)\left(\sum_{i=0}^{N}x_{i}\right)}{\left(\sum_{i=0}^{N}x_{i}^{2}-\frac{1}{N}\left(\sum_{i=0}^{N}x_{i}\right)^{2}\right)}$$

Ou seja:

$$b_{1}=\frac{\sum_{i=0}^{N}x_{i}y_{i}-N\overline{y}\overline{x}}{\sum_{i=0}^{N}x_{i}^{2}-N\overline{x}^{2}},\qquad b_{0}=\overline{y}-b_{1}\overline{x}$$
A partir dessas constantes, montamos a seguinte equação da reta:
$$y=b_{0}+b_{1}x$$
Comparando com a lei de Pareto linearizada, temos que a constante de regressão  $b_0$ é  $b_0=b$,  e o coeficiente de regressão $b_1$ é o próprio índice Pareto $\alpha=b_1$  que estávamos buscando.  Devemos lembrar que obtemos a forma linearizada da lei de Pareto apenas após aplicarmos $\log_{10}$. Dessa forma, para verificar se um conjunto de dados segue uma distribuição de Pareto, devemos aplicar a regressão linear somente depois de calcular a CCDF dos dados e, em seguida, aplicar $\log_{10}$ em ambos os eixos. Em outras palavras, a verificação é feita no gráfico $\log_{10} (w) \times \log_{10} (\overline{F}_x(w))$.


\subsubsection*{Gini}

%{\color{red} Talvez escrever sobre o gini?}

%https://www.econstor.eu/bitstream/10419/242579/1/aut-ewp202007.pdf

A curva de Lorenz é dada por uma função do tipo $L(x)$ que nos indica a fração de riqueza (ou renda) em posse da fração $x$ da população. A igualdade perfeita é dada, então, por uma reta $L(x)=x$.

O coeficiente de Gini é uma medida da desigualdade em um conjunto de dados e é dado pela razão entre a área que se encontra entre a reta de igualdade perfeita e a curva de Lorenz, e a área total sob a reta de igualdade perfeita.

Se temos um conjunto de dados ordenados $(y_1\leq y_2\leq \dots\leq y_N) $, a proporção acumulada da população, ou seja, a posição no eixo $X$ para o dado $i$ é dada por:

$$x_i=\frac{i}{N}$$

Consequentemente a distância entre dois dados $i$ e $i+1$ é dada por

$$\Delta x =\frac{1}{N}$$

A altura da curva de Lorentz é dado por:

$$y_i=\frac{\sum_{k=1}^{i}y_{k}}{\sum_{k=1}^{N}y_{k}}$$

Para a reta de igualdade perfeita temos:
$$y_{i}^{*}=x_{i}=\frac{i}{N}$$

Logo a diferença entre as alturas é dada por:

$$\Delta y_i=\frac{i}{N}-\frac{\sum_{k=1}^{i}y_{k}}{\sum_{k=1}^{N}y_{k}}$$

Assim sendo, a área dada pela diferença entre a reta da igualdade e a curva de Lorenz para a distância entre dois dados é dada por:

$$\Delta x\Delta y_i=\frac{1}{N}\left(\frac{i}{N}-\frac{\sum_{k=1}^{i}y_{k}}{\sum_{k=1}^{N}y_{k}}\right)$$

Somando então sobre todos os dados:

\begin{equation}
    \begin{split}
        A	&=\sum_{i=1}^{N}\frac{1}{N}\left(\frac{i}{N}-\frac{\sum_{k=1}^{i}y_{k}}{\sum_{k=1}^{N}y_{k}}\right) \\
	&=\frac{1}{N}\left(\frac{1}{N}\sum_{i=1}^{N}i-\frac{\sum_{i=1}^{N}\sum_{k=1}^{i}y_{k}}{\sum_{k=1}^{N}y_{k}}\right)
    \end{split}
\end{equation}
O primeiro termo entre parênteses:
$$\frac{1}{N}\sum_{i=1}^{N}i=\frac{1}{N}\frac{N\left(N+1\right)}{2}=\frac{N+1}{2}$$
Para o somatório duplo no segundo termo, percebendo que para cada $k$  o valor $y_k$ vao aparecer em todos $i\geq k$, isto é $N-k+1$ vezes:
$$\sum_{i=1}^{N}\sum_{k=1}^{i}y_{k}=\sum_{k=1}^{N}\left(N-k+1\right)y_{k}$$
Então:
$$A	=\frac{1}{N}\left(\frac{N+1}{2}-\frac{\sum_{k=1}^{N}\left(N-k+1\right)y_{k}}{\sum_{k=1}^{N}y_{k}}\right)$$
Por fim, nos resta apenas lembrar que o Gini é dado pela razão entre a área que se encontra entre a reta da igualdade e a curva de Lorenz ($A$) e a própria área abaixo da reta da igualdade. Considerando que os dados estejam normalizados, a área abaixo da reta da igualdade perfeita corresponde simplesmente à metade da área de um quadrado de lado $1$ ou seja, $1/2$ o coeficiente de Gini é então:

$$G	=\frac{2}{N}\left(\frac{N+1}{2}-\frac{\sum_{k=1}^{N}\left(N-k+1\right)y_{k}}{\sum_{k=1}^{N}y_{k}}\right)$$
Ou em uma forma mais comumente encontrada na literatura\cite{gini}:
$$G=\frac{1}{N}\left(N+1-2\left(\frac{\sum_{k=1}^{N}\left(N-k+1\right)y_{k}}{\sum_{k=1}^{N}y_{k}}\right)\right)$$